\section{Goals and Timetable}
\label{sec:goals}

\subsection{Original goal}

The original goal was \textbf{efficiency}: make the paper's detector cheaper to run without
changing its answers. The idea was an \emph{edit classifier}. It would look at the query
image first, predict what kind of edit had been applied, and run only the BK-tree searches
that could still contribute a vote. For example, a pixelated image destroys the fine detail
that a structure-based signal needs, so that signal's search could be skipped.

The planned steps were:
\begin{enumerate}
  \item reimplement the paper's four hashes in C11, bit-exact against the reference library;
  \item build a labelled dataset of NFT copies with known edits;
  \item train the edit classifier;
  \item benchmark the classifier-guided detector against the original for speed.
\end{enumerate}

\subsection{How the goals changed}

To build the classifier we had to understand exactly where each signal breaks, and that
work changed the project in three ways:
\begin{itemize}
  \item \textbf{sHash cannot be searched correctly} (Section~\ref{sec:shash}). Accelerating
        four BK-tree searches makes little sense when one of the four trees can miss real copies.
  \item \textbf{A C11-versus-Python speed comparison would not be a fair claim.} Our
        instructor pointed out that comparing a C11 implementation against the paper's Python
        mixes language effects with algorithmic ones.
  \item \textbf{A better target existed.} The paper's own future-work section points at its
        geometric blind spots, and feature matching addresses exactly those.
\end{itemize}

The revised goals were therefore about \textbf{accuracy}:
\begin{enumerate}
  \item replace sHash with a signal that does its job, cropping and rotation, and can be indexed
        soundly;
  \item measure the replacement fairly: against sHash on its own edits, inside the paper's
        unchanged detector, and against its storage cost;
  \item re-aim the edit classifier from speed to accuracy (stop trusting a signal on the edits
        known to break it) and measure whether that helps.
\end{enumerate}
The C11 hash suite was kept as validated groundwork. The rest of the project was done in
Python.

\subsection{Timetable}

\begin{table}[ht]
\centering\small
\begin{tabularx}{\linewidth}{@{}l X X@{}}
\toprule
Stage & Work & Outcome \\
\midrule
1. Learning & The paper, perceptual hashing, BK-trees, the \texttt{imagehash} code & Project plan \\
2. Reproduction & C11 hash suite; differential tests on 600 images and the authors' data &
  Bit-exact hashes; the sHash finding \\
3. Pivot & sHash analysis; review of a first ORB pipeline & Old ORB numbers retired; pairwise
  evaluation adopted \\
4. Phase B: ORB & Pairwise ORB signal, threshold tuning, comparison with sHash, storage curve &
  Results 1--2; storage finding \\
5. Phase C: edit classifier & 93 features, Random Forest, per-signal reliability map &
  Accurate classifier \\
6. Phase D: full detector & Three-way comparison, comparison at the paper's own thresholds,
  cross-generator check & Result 3; classifier adds $\approx 0$ \\
7. Phase E: multi-edit test & Two-edit generator, independent classifier heads, four-way
  comparison & The colour trace is erased by pixelation \\
8. Write-up & Presentation, final report, poster & This document \\
\bottomrule
\end{tabularx}
\caption{Project stages, in order.}
\label{tab:timetable}
\end{table}
